\section{Conclusion}


In this paper, we propose a unified dual-mask physical model for light field depth estimation to address the severe performance degradation caused by non-Lambertian reflectance and complex textures. By introducing a physical-level angular signal decoupling mechanism and an adaptive dual-mask routing architecture, our framework effectively mitigates the conflict between Epipolar Plane Image (EPI) geometric assumptions and complex optical phenomena, achieving superior accuracy in mixed and overall scenes.

The core contributions of this work are summarized as follows:
\textit{ \textbf{Three-Layer Angular Signal Decomposition and Geometric Material Classification:} We introduced a physical-level decoupling mechanism that explicitly separates illumination, depth, and material properties. This resolves spectral feature degradation under low angular resolutions (e.g., $9 \times 9$) and provides robust material priors, enabling the accurate identification of non-Lambertian regions that inherently violate EPI assumptions.
} \textbf{Geometry Dual-Mask Routing for Unified Depth Estimation:} We constructed an adaptive routing architecture that overcomes the bottleneck of single-EPI networks in handling mixed reflectance properties. By dynamically processing complex reflections within a unified model, our method significantly enhances depth accuracy in mixed and overall scenarios, achieving an Overall MAE of 0.133 and a Mixed (Urban) MAE of 0.081.
\textbf{Quantitative Analysis of EPI Assumption Failure Boundaries:} We rigorously defined the theoretical performance boundaries and root causes of EPI assumption breakdowns. By quantifying the inherent architectural limitations of EPI-based networks through 132 extensive empirical trials, this analysis prevents ineffective computational efforts and provides a solid theoretical foundation for future non-EPI paradigms and non-Lambertian dataset construction.

Finally, since our extensive empirical analysis demonstrates that single EPI architectures possess an inherent physical ceiling when resolving extreme non-Lambertian reflections and texture-induced slope ambiguities, future work will pivot towards non-EPI paradigms. Specifically, we plan to integrate neural radiance fields (NeRF) and learned feature matching techniques to transcend the theoretical boundaries of epipolar geometry and achieve universally robust light field depth estimation.


From a signal processing perspective, the superiority of the proposed three-layer angular signal decomposition over traditional frequency-domain methods lies in its resilience to angular undersampling. Conventional approaches often rely on 2D Discrete Fourier Transform (DFT) to extract angular frequency spectrums. However, under low angular resolutions (e.g., $9 \times 9$), the angular sampling rate frequently falls below the Nyquist limit for complex textures and sharp specularities. This induces severe spectral leakage and aliasing, entangling depth-induced low-frequency shifts with material-induced high-frequency components. By formulating the light field signal as $I(u,v) = DC + a\cdot u + b\cdot v + \epsilon(u,v)$ in the spatial domain, our model inherently acts as a robust parametric filter. The linear terms $(a, b)$ capture the low-frequency epipolar geometry, while the residual $\epsilon(u,v)$ isolates the high-frequency non-Lambertian perturbations. This spatial-domain regression not only circumvents the spectral degradation inherent in DFT-based pipelines but also significantly reduces the computational overhead associated with high-dimensional frequency transformations.

Furthermore, the explicit extraction of geometric priors from the residual term $\epsilon(u,v)$ provides a distinct advantage in cross-domain generalization compared to implicit end-to-end deep learning baselines. Convolutional neural networks trained for material classification often overfit to dataset-specific texture biases and absolute illumination conditions. In contrast, our physical residuals are intrinsically invariant to global illumination shifts. Features derived from $\epsilon(u,v)$, such as the angular gradient and cone score, capture the fundamental photometric inconsistencies caused by specular lobes and volumetric scattering. The exceptionally high separability of these features ensures that the Random Forest classifier operates on a highly discriminative, physically grounded manifold. Importantly, utilizing a lightweight, non-differentiable classifier for this phase prevents the material prior extraction from burdening the main GPU memory during end-to-end depth regression, enabling efficient deployment while preventing the network from relying on spurious correlations in unseen urban environments.

The introduction of the geometric dual-mask routing architecture also resolves a critical optimization bottleneck in unified depth estimation: gradient interference. In standard single-branch networks, the simultaneous processing of Lambertian and non-Lambertian pixels creates conflicting gradient directions during backpropagation. The epipolar consistency loss pulls the network toward linear EPI structures, while complex reflections demand non-linear feature adaptations. This conflict flattens the loss landscape and traps the model in suboptimal local minima. By dynamically routing pixels based on the physical material priors, our framework effectively decouples these gradient flows. Coupled with domain-balanced sampling, this mechanism allows the network to apply specialized physical constraints to distinct reflectance domains without catastrophic forgetting. This smooths the optimization trajectory, accelerates convergence, and yields significantly more precise depth boundaries in mixed-reflectance scenes where traditional EPI assumptions collapse.

Despite these advancements, the current framework exhibits certain limitations rooted in its underlying physical assumptions. The linear parameterization of the depth term assumes local planarity and unoccluded visibility across the micro-lens array. Consequently, the model may experience performance degradation at severe occlusion boundaries or when imaging highly curved transparent and refractive surfaces, where light transport becomes inherently non-linear and viewpoint-dependent. Additionally, the current formulation processes individual light field frames independently, ignoring temporal redundancies and multi-view geometric constraints over time. Future work will explore integrating non-linear light transport models for refractive materials and extending the dual-mask routing mechanism to spatiotemporal light field video sequences. Incorporating active illumination cues or multi-modal sensor fusion could further enhance depth recovery in extremely low-light or highly absorptive non-Lambertian scenarios.

Despite the demonstrated efficacy of the proposed unified dual-mask physical model, several limitations remain, primarily stemming from the underlying signal processing assumptions and computational constraints. First, the three-layer angular signal decomposition relies on a local linear approximation, formulated as $I(u,v) = DC + a\cdot u + b\cdot v + \varepsilon(u,v)$. While this effectively decouples illumination and depth for standard Lambertian and mildly specular surfaces, it struggles with highly anisotropic materials, multi-layered translucent media, or severe occlusions. In such extreme scenarios, the angular signal exhibits strong non-linearities and abrupt discontinuities, causing the residual term $\varepsilon(u,v)$ to inadvertently absorb geometric depth cues. This spectral leakage can lead to false positives in the geometric material classification, subsequently misrouting pixels within the dual-mask architecture.

Furthermore, the current implementation utilizes a Random Forest (RF) classifier operating on explicitly extracted geometric features. Although this provides high interpretability and robust performance under low angular resolutions (e.g., $9 \times 9$), shallow machine learning models inherently lack the representational capacity to generalize to out-of-distribution (OOD) illumination conditions or unseen complex Bidirectional Reflectance Distribution Functions (BRDFs). Moreover, the reliance on explicit feature extraction prior to classification creates a processing bottleneck. Additionally, the pixel-wise physical decomposition and the adaptive dual-mask routing mechanism introduce non-trivial computational overhead, which currently restricts the deployment of the framework in latency-critical embedded systems, such as real-time autonomous driving platforms.

To address these limitations, future work will explore several promising directions. We plan to extend the linear decomposition model by incorporating higher-order polynomial expansions or microfacet-based non-linear formulations to better capture anisotropic reflections and complex occlusion boundaries. Furthermore, integrating self-supervised learning paradigms could alleviate the dependency on precise material annotations. Instead of relying on a decoupled RF classifier, we aim to develop a differentiable, physics-informed neural module that allows for joint end-to-end optimization of material priors and depth estimation, thereby enhancing generalization to OOD scenarios. Finally, extending the current spatial dual-mask routing into the spatiotemporal domain represents a critical next step. By leveraging temporal consistency constraints across consecutive light field video frames, we can further regularize the depth estimation in highly dynamic non-Lambertian regions, ultimately bridging the gap between theoretical light field models and real-world robotic manipulation applications.